\chapter[Data-driven estimation of symbolic representations of egocentric vision for the assessment of dysexecutive syndromes]{Data-driven Estimation of Symbolic Representations of Egocentric Vision for the Assessment of Dysexecutive Syndromes}
\label{sec:chapter_5}


\begin{chapterabstract}
    \textbf{Key points:}
    \begin{enumerate}
        \item The behavioral alterations induced by dysexecutive syndromes (DS) during a complex prescribed task are diverse, heterogeneous, and manifest spontaneously, posing challenges to their identification and measurement. This encourages the development of novel computational methods to encode participants' behavior more systematically at fine granularity.

        % \item The nascent field of computational ethology aims to develop automatic, objective, and scalable methods to quantify animal (including human) behavior in naturalistic settings. Behavior is hypothesized to be built from sequences of stereotyped modules of actions, often organized hierarchically within ethograms. 

        \item Fine-grained and temporally accurate symbolic representations of participants' egocentric vision offer rich, compact, and interpretable proxy representation of behavior to study task execution, amenable to pattern mining and discovery, and facilitating multimodal integration.

        \item  We leverage the redundancy and temporal regularities of egocentric vision in repeatable settings to construct a generalizable vocabulary of prototypes. Prevalent and unequivocal atomic actions are discovered by iteratively refining latent subspaces using minimalistic human feedbacks to ascertain their quality and recover their interpretability. Prototypes are then defined as sub-trees of a silhouette score-optimized hierarchical model to reduce redundancy and lower the vocabulary cardinality. 

        \item Our symbolization combine the prototype vocabulary with a fine-grained temporal segmentation of participants' visual trajectory that solve a multiple change-point detection problem exactly and efficiently.

        %enabling the use of machine learning and pattern mining technics while facilitating multimodal intergration that could enhance the study of behavioral manifestation of DS.  %COde: move to the key points ! 

        % \item These symbolic representations provide interpretable and scalable tools for assessing automatically established markers of executive dysfunction, offering compact yet expressive surrogate for behavioral pattern mining and discovery, and facilitate multimodal integration.      

        % \item Clustering and temporal regularization are combined to estimate robust and reliable egocentric vision ethograms, which are symbolic representations of the participant's visual trajectory during the cooking task.

    \end{enumerate}
    \textbf{Contributions:}
    \begin{enumerate}
        %  \item We propose an iterative exploration of carefully defined clustering residual subspaces of empirical egocentric data data-driven computational method with human feedbacks to establish a visual phenotype of participant's egocentric vision during the cooking task.

         \item We propose a novel computational framework to encode automatically and accurately participants' egocentric vision during a cooking task as fined-grained symbolic representations. The resulting symbolic sequences allow to study differences in task execution across diagnosis groups to support the assessment of dysexecutive syndromes.

        \item We demonstrate the generalizability and accuracy of the symbolic representations in unsupervised settings via powerful visual probing mechanisms, and evaluate the robustness to human feedback subjectivity through the construction of two independent models yielding similar vocabularies.
        
        
    \end{enumerate}
\end{chapterabstract}



\section{Introduction} %Difference with action rcognitione etc 
\label{sec:chapter_5_introduction}

    Recent advances in machine learning (ML) and computer vision (CV) offer novel opportunities for the characterization of intricate and multifaceted behavioral alterations induced by neurological and psychiatric disorders. Systematic recording of participants behavior during a complex ecological task offer unprecedented opportunity to characterize these manifestations more accurately, holding promises to understand the links between brain lesions, neurological impairments, and behavioral alterations. However, their spontaneous manifestation through semantically rich and sequentially complex and diverse actions poses a significant challenges for computational methods to identify them automatically in ecologically relevant settings. 

    The nascent field of computational ethology aims to fill this gap by developing automatic, objective, and scalable methods to quantify animal (including human) behavior accurately in naturalistic settings, in all its richness and complexity \cite{Anderson2014, Wiltschko2015, Berman2014}. For psychiatric and neurological conditions, it has the potential to automate and objectify the identification of complex symptoms, and advance continuous, longitudinal patients monitoring and treatment approaches \cite{Bavaresco2022, Mobbs2021}. 

    Since its inception, a central hypothesis in ethology formulates that behavior is built from an identifiable sequence of stereotyped modules of actions \cite{tinbergen1951study}. These sequences of actions were suggested to follow central plans that exhibit a hierarchical structure \cite{BOTVINICK2008201, Rosenbaum2007, Barker1965, lashley1951serial}: at the finest scale, behavior consists of atomic body movements (for example "hand grasping" or "head turning"), while at coarser levels, actions are more naturally defined by goal-directed behaviors and intentionality (e.g. "adding the eggs" or "cook the dough"). In ethology, naturalistic behaviors are often standardized using ethograms, which are structured catalogs of behaviors relevant to a specific research question \cite{tinbergen1951study}. Ideally, ethograms aim to partition behavior into \textit{mutually exclusive} and \textit{exhaustive} categories, in order to provide a complete reference framework for systematic observation of behavior \cite{Jones2016Etho}. For example, \Cref{fig:chapter_4_hierarchical_behavior} illustrates the decomposition of a simple coffee-making task into increasingly fine-grained behavioral units \cite{Cooper2000}.
    
    \begin{figure}[!htbp]
        \centering
        \includegraphics[width=0.8\linewidth]{chapters/chapter 4/figures/hierarchy_bis.png}
        \caption{Hierarchical organization of a coffee-making task, adapted from \cite{Cooper2000}.}
        \label{fig:chapter_4_hierarchical_behavior}
    \end{figure}

    Traditional ethograms, while valuable, face major limitations: (i) subjectively pre-defining action categories; (ii) time-consuming manual annotations; (iii) inter-rater reliability; and (iv) the number, granularity, and temporal resolution of categories are constrained. These issues render ethograms difficult to scale to complex and long tasks such as cooking an entire recipe. In contrast, unsupervised computational methods aim to automatically identify and extract stereotyped behavioral modules (also called primitives, motifs, or "syllabes") from large-scale recordings, holding the promise for more objective, scalable, and accurate quantification of behavior \cite{Berman2014, Wiltschko2015, Datta2019}.
    
    These approaches typically start by extracting low-dimensional body representations (e.g. pose estimation, kinematics data), before clustering these features into temporally consistent actions by detecting recurrent and stereotyped patterns \cite{Vogg2025, Zhang2024}. The hierarchy of actions is often reconstructed using trees, where the leaves correspond to the observable atomic behaviors, and higher-level nodes encode more complex actions formed by grouping inferior nodes based on their transition patterns \cite{Miller2024, HOP_truong}. Alternatively, supervised annotation can be used to recover higher-level actions semantics, by mapping atomic actions to higher-level categories.

    However, the recent advance in animal computational ethology \cite{FAZZARI2025128330, Tang2024}, driven by the maturation of deep learning-based automatic pose estimation (e.g. DeepLabCut \cite{Mathis2018}, SLEAP \cite{Pereira2022}, DeepPoseKit \cite{Graving2019}, Deep Graph Pose \cite{DeepGraphPose}) and temporal segmentation algorithms (e.g. KeyPoint-MoSeq \cite{Weinreb2024}, SimBA \cite{Goodwin2024}, VAME \cite{Luxem2022}), does not easily transfer to human behavioral analysis, in particular for detecting subtle behavioral variations during complex ecological tasks \cite{Datta2019}. 
    
    %Other work \cite{ARIFOGLU201786} While some work study behavioral alterations as they degrade over time using, using longitudinal data 
    In Parkinson disease for example, wrist-worn wearable kinematics were used to generate symbolic representations of movements that successfully characterized severity-stratified motor disruption during in-clinics and at home activities \cite{Abrami2020}. Using video data, a recent study combined pose tracking with unsupervised behavioral atoms segmentation to quantify the temporal structure, transition dynamics, expression variability, and stereotypy of bipolar disorder during a free-moving exploratory task \cite{Zhang2024}. However, the majority of studies on human and nonhuman behavior still rely primarily on body motion analysis, which miss the integration of multimodal and contextual cues (e.g. objects, social interactions, pupillometry and oculometry) that encompass behavior. In particular, the identification of participant's progression through the task is a strong requirement of our behavioral assessment method, ensuring the interpretability and feasibility of downstream analysis, and facilitating multimodal integration through context incorporation.%, which mosly arise from participant's environment cues captured by the videos.  
    
    % We belive the complexity regime of egocentric vision atoms is different from motion-based atoms, requiring specific methodological developments. (in practice we need to test the VAME framework to our setting)

    Egocentric sensing technologies, such as eye-trackers and head-mounted cameras, provide invaluable sources of information for advancing behavioral analysis \cite{RAGUSA2023103764, NUNEZMARCOS2022175}. In contrast to exocentric recordings, they capture fine-grained and participant-centered contextual information. This includes pupillometry and oculomotor dynamics, which provide access to physiological and cognitive states \cite{Castellotti2025, Strauch2022, Walter2021}, high-resolution and occlusion-free recordings of hand–object interactions, first-person visual perspectives that disambiguate actions, and integrated cues of gaze, head, and body movement \cite{NUNEZMARCOS2022175}. Together, these signals constitute a uniquely detailed account of how participants perceive and interact with their environment \cite{7298625}.

    In this work, we propose to encode the participants visual experience during the task into symbolic representations. Specifically, we model visual experiences as sequences of fine-grained visual atoms, where each atom denotes an interpretable and unequivocal action. Our hypothesis is that the redundancy and temporal regularity of short segments of egocentric vision during a standardized task can be exploited to build a common vocabulary of vision prototypes. To the best of our knowledge, this is the first time an ethologically relevant vision-based behavioral phenotyping approach is proposed for the analysis of complex human behaviors during ecological tasks.
    % Add a few sentences about the importance of the task structure to guide the analysis.
    \paragraph{Chapter outline.\\}
    
    We first present the three components of our symbolic representation method, which consists in (1) an egocentric vision prototyping stage (\Cref{sec:chapter_5_prototyping}), a (2) temporal segmentation module (\Cref{sec:chapter_5_unsupervised_temporal_segmentation}), and (3) the symbolic representation computations (\Cref{sec:chapter_5_symbolic_representation}). Then we present our experimental setting in \Cref{sec:experimental_setting}, and unfold the evaluation protocol with associated results in \Cref{sec:chapter_5_evaluation_protocol}. Finally, we conclude in \Cref{sec:chapter_5_conclusion}.

\section{Egocentric vision prototyping using recursive latent subspace exploration with human feedbacks}
    \label{sec:chapter_5_prototyping}

    \subsection{Definitions and problem formulation}

    The following definitions will be used in this chapter and the following one:

    \begin{itemize}
        \item A \textbf{time series} is a sequence of real-valued vectors, denoted as $\mathbf{X} = (\mathbf{x}_1,\ldots, \mathbf{x}_n)$, where $\mathbf{x}_i \in \mathbb{R}^D$ is the feature vector at index $i$, $D$ is the feature dimension, and $n$ is the length of the time series. 

        \item A \textbf{string} is a sequence of discrete values indexed in time order. The alphabet of non-ordered symbols is denoted $\mathcal{A} = \{a_1, \ldots, a_G\}$, of cardinality noted $G=|\mathcal{A}|$.

        \item A \textbf{symbolic sequence}, also called symbolic representation in this work, is a discrete sequence that encodes a multivariate time series, denoted as $\mathbf{s} = (s_1, \ldots, s_n)$, where $s_i$ is the symbol associated with the feature vector $\mathbf{x}_i$ at index $i$. The symbolization procedure may induce an order over the alphabet of symbols, which is discussed in \Cref{sec:chapter_6}.

        \item An \textbf{prototype} of egocentric vision is an interpretable and unequivocal visual pattern that represents a class of semantically consistent visual experience. We denote $\mathcal{V}$ a vocabulary of prototypes, using their (non-ordered) indexes to define an alphabet of symbols. We sometimes use the term "visual atom" to refer to a prototype, and omit "egocentric vision" qualifier for conciseness.

        % \item A \textbf{visual trajectory} $\mathbf{X}^(k)$ is the multivariate time series of feature vectors extracted with the VideoMAE-v2 encoder from the egocentric video of participant $k$ during the cooking task. The length of the visual trajectory, denoted as $N_i$, may vary between participants.

    \end{itemize}

    The goal of this section is to construct a representative vocabulary $\mathcal{V}$ of egocentric vision prototypes, generalizable across participants, and sparse, that is, of minimal cardinality. Akin to visual phenotyping, the vocabulary should be sufficiently rich to capture the diversity of visual experiences encountered during the task, while being compact enough to facilitate interpretability and downstream analysis. In the vision domain, we relax the exhaustivity constraint of ethograms in favor of statistical coverage, defining our goal as the identification of prototypical visual atoms, which is akin to partitioning the VideoMAE-v2 latent space to define semantically unequivocal regions. % Finite vocab and mutual-exclusivity removed ..

    To illustrate the underlying clustering challenge and motivate our approach on real data, we define naive prototypes as the centroids estimated using cosine $k$-means clustering on all available videoMAE-v2 embeddings for increasing vocabulary size $G$. We visualize the resulting symbolic representations of participants in \Cref{fig:chapter_5_assignments_across_K}. 

    \begin{figure}[!htbp]
        \centering
        \includegraphics[width = 1.0\linewidth]{chapters/chapter 5/figures/chapter_5_assignements_across_K.png}    
        \caption{\textbf{Participant's symbolic representation for increasing vocabulary size.} Each row corresponds to a different participant performing the cooking task, and each color represents a different cluster. Sequences are upsampled to match the longest sequence for visualization purposes. Coarse clusters reveal inter-subjects temporal consistency, while finer clusters capture more specific visual atoms.}
        \label{fig:chapter_5_assignments_across_K}
    \end{figure}

    We first observe a remarkable temporal consistency of prototype assignments across participants, especially for low $G$, suggesting the feasibility of learning a generalizable vocabulary, and raising the question of the minimal vocabulary size required to capture the diversity of visual experiences. In order to relate these clusters to the structure of participant's similarity matrices and clarify clustering behavior, we compute and visualize in \Cref{fig:chapter_4_gram_across_participants} the cosine similarity matrices of six participants from the control and patient groups. 

    \begin{figure}[!htbp]
        \centering
        \includegraphics[width = 1.0\linewidth]{chapters/chapter 5/figures/gram_participant_raw.png}
        \caption{\textbf{Pairwise cosine similarity matrix between time-aligned embeddings for three control participants (top) and three patients with radiation-induced leukoencephalopathy (RIL) (bottom).} The color scale represents the cosine similarity between pairs of latent embeddings, with warmer colors indicating higher similarity. These similarity matrices provides rich information about participant's execution of the task, and illustrate the efficient representational power of the videoMAE-v2 video encoder. For example, key steps can be identified, such as reading the recipe (purple bar on the left), breaking the eggs (orange bar), as highlighted in the figure. Our prototyping approach focus on the smallest diagonal blocks of stable elevated similarity. %Notable events, such as reading the recipe and breaking the eggs, are highlighted. CODE which 
        }
        \label{fig:chapter_4_gram_across_participants}
    \end{figure}

    The organization of participants’ visual trajectory combines (i) short sequential atomic actions at fine granularity, and (ii) hierarchical patterns emerging within larger diagonal blocks, combining multiple visual atoms. While this multi-scale organization of visual experiences is reminiscent of the multi-scale structure of human activity discussed above, where fine-grained atomic actions are nested within higher-level goal-directed behaviors, these may only partially align since cosine similarity primarily reflects visual similarity rather than semantic task boundaries.
    
    These experiments suggest that coarse prototypes (low $G$) correspond to large-scale diagonal blocks merging several atomic visual atoms, while finer-grained prototypes (high $G$) more likely capture specific and unequivocal visual atoms. Within this multi-scale landscape of pairwise similarities, our goal is to focus on the finest diagonal blocks, where egocentric vision segments primarily reflects stable short-term motion and appearance. %While this scale allows a simpler and more objective definition of atomic units of action, we aim to address the challenge of constructing a fine-grained vocabulary of \textit{unequivocal} visual atoms that can accurately reflect participants visual experience. 
    However, in our high-dimensional ($D=1408$) and high-complexity setting (e.g., $K=100$), we observed that a direct clustering of the empirical sample ($N \approx 906k$) - sometimes referred to as extreme clustering setting \cite{10.1145/3097983.3098079} — does not readily produce a reliable vocabulary, due to clusters emerging at variable scales. Indeed, using a methodology we describe in \Cref{sec:chapter_5_annotation} to recover clusters semantics, we observe that, although most estimated clusters capture behaviorally interpretable and unequivocal actions, a subset remains ambiguous and mix multiple visual atoms. This limitation holds across clustering methods, dimensionality reduction, and normalization strategies. 
    
    Building on these observations, we propose a novel method to partition the latent space recursively by (i) accumulating the fined-grained unequivocal prototypes, and (ii) refining under-segmented regions associated with ambiguous clusters. This recursive latent subspaces exploration is guided by minimalistic human feedbacks to ascertain clusters quality and recover their interpretability. 
     
    \subsection{Description of the prototyping procedure}
    \label{sec:chapter_5_recursive_procedure}

        Our prototyping procedure estimates a vocabulary in two stages. The first stage, described in this section, consists in estimating an initial set of $K$ initial centroids through the iterative exploratory procedure incorporating human feedback. The second stage, detailed in \Cref{sec:chapter_5_hierarchical_multiview_clustering}, consolidates the set of accumulated centroids to favor model sparsity by gathering semantically redundant clusters via hierarchical clustering. %This refinement step strives to gather \textit{redundant} visual atoms that belong to the same semantic category, by incorporating centroids temporal structure using a kernel-based distance-level fusion. Specifically, we balance the centroids pairwise distance in the representation space with the Wasserstein distance between their temporal distributions, improving the accuracy of the consolidation stage. 

        \paragraph{Centroids accumulation step.\\}

        Let $X_{\text{train}} = \{ x_1, ..., x_N \} \in (\mathbb{R}^{D})^N$ denote a set of $N$ training latent embeddings. The first centroid accumulation stage proceeds iteratively over $P$ steps. At each step $p \leq 1$, the following two steps are performed:

        \begin{enumerate}
            \item \textbf{Generative step.} 
            
            Estimate $C$ candidate centroids $\mu^{p} = (\mu^p_1, ..., \mu^p_C) \in \mathbb{R}^D$ by applying cosine $k$-means clustering to the residual set $R_p$ (defined below). The centroids are estimated by solving the optimization problem
            \begin{equation}
                \hat{\mu}^p = \arg \min_{\mu^p_{1},\dots,\mu^p_{C}} \;
                \sum_{x_i \in R_p} \Bigl(1 - \tfrac{x_i^\top \mu^p_{c_i}}{\|x_i\|\;\|\mu^p_{c_i}\|}\Bigr),
            \label{eq:chapter_5_kmeans_objective}
            \end{equation}

            where $c_i \in [C]$ is the cluster assignment of $x_i$, defined as $c_i = \arg \min_{c \in [C]} (\tfrac{x_i^\top \mu^p_{c}}{\|x_i\|\;\|\mu^p_{c}\|})$

            \item \textbf{Discriminative step.} 
            
            Each of the $C$ empirical clusters are manually reviewed using the annotation protocol defined in \Cref{sec:chapter_5_annotation}, to separate the centroids indexes into two disjoint subsets: 
            
            \begin{itemize}
                \item $I_p^+ \subset [C]$, associated with fine-grained and unequivocal clusters, and 
                \item $I_p^- = [C] \setminus I_p^+$, associated with semantically ambiguous clusters. 
            \end{itemize}
            We denote by $A_p$ the set of accumulated centroids up to step $p$, defined as $A_p = \bigcup_{k=1}^p \{ \mu^k_j\}_{j\in I_k^+}$, and $N_p = \{ \mu^p_c \}_{c \in I_p^-}$ the set of rejected centroids.

            \item \textbf{Samples assignments update.} Samples assignment is updated using centroids from $A_p$ and $N_p$ to define the residual step at step $p+1$. We denote $d_i = \min_{c \in A_p \cup N_p} d(x_i, \mu_c)$ the assignment distance of sample $x_i$ to its closest centroid.

            \item \textbf{Residual set update.} The residual set, $R_{p+1}$, is defined from two sources of empirical samples: 
            \begin{itemize}
                \item Samples associated with the ambiguous subset of centroids $N_p$.
                \item Samples associated with valid centroids in $A_p$ whose assignment distance is above a centroid-specific threshold, noted $d^{\text{min}}_c$ for centroid $c \in A_p$.
            \end{itemize}
            The initial residual set comprises all samples, that is, $R_1 = X_{\text{train}}$.

        \end{enumerate}     

        The procedure ends at step $P$ when the accumulated number of source prototypes is above a predefined number $K_{\text{max}}$, yielding a set of $K$ prototypes, $\hat{\mu}^{[K]} = A_P \hat{=}  (\hat{\mu}_1, ..., \hat{\mu}_K)$.

        The residual set is constructed to emphasize regions of the latent space that are misrepresented by current cumulative set $A_p$, corresponding to under-segmented regions. By isolating them in subsequent clustering steps, we aim to refine clusters' semantic and encourage the discovery of semantically novel clusters. The procedure is illustrated in \Cref{fig:chapter_5_recursive_prototypes_workflow}.
        
        \begin{figure}[!htbp]
            \centering
            \includegraphics[width = 1.0\linewidth]{chapters/chapter 5/figures/chapter_5_recursive_prototypes_workflow.png}    
            \caption{\textbf{Recursive latent space exploration and prototyping workflow.} The workflow iteratively explores under-segmented regions of the videoMAE-v2 latent space to extract fine-grained cluster centroids guided by human feedback. It starts from an initial set of centroids obtained with K-means ($K=100$ in our experiments), which are then evaluated and refined across successive iterations with human annotators. This recursive process facilitates the discovery of unequivocal and representative centroids that capture core patterns of participants' egocentric videos.}
            \label{fig:chapter_5_recursive_prototypes_workflow}
        \end{figure}

        \begin{rmk}
            We used cosine $k$-means efficiently implemented in the \texttt{Faiss} library, as it was found in a preliminary experiment to outperform other implementation and distance metrics in terms of number of retained clusters in the first round, that is, maximizing $|I_1^+|$, which was used as a proxy to quantify clustering algorithms quality. We report this experiment in  \Cref{apd:chapter_5_clustering_technics}. % CODE: how about the SSDE for each of these solution ? They are not comparable but maybe using distribution comparison (more left skewed is etter)
        \end{rmk}

        \begin{rmk}
            The condition on the assignment distance in the definition of the residual space is discussed in \Cref{apd:chapter_5_distance_distribution_thresholding}. In the experiments, we use a constant value $d_c^{\text{min}}=0.3$ for each accepted centroid $c$, (close to the first quartile of assignment distances values using $A_1$), in favor of a stringent constraint on the assignment distance. 
            % \item We stop the procedure after the cumulated number of centroids is significantly higher than the average number of visual segments detected when solving the UTS task, which is discussed in \Cref{apd:chapter_5_stopping_criteria}.
        \end{rmk}       
        
        In the following, we omit the hat notation for simplicity, and denote the final set of $K$ source centroids as $\mu^{[K]} = (\mu_1, ..., \mu_K)$.



    \subsection{Annotation protocol for the egocentric vision prototypes}
        \label{sec:chapter_5_annotation}

        Recovering the clusters interpretability is crucial for ensuring meaningful interpretations of the symbolic representations in downstream analysis tasks. In practice, however, coarse clusters may emerge at scales where participants' visual experiences encompass a broad range of action semantics, as illustrated in \Cref{fig:chapter_5_assignments_across_K}. Such clusters hinder clear interpretation and prevent the definition of unequivocal, fine-grained motifs that characterize a unique visual experience.

        To address this, we introduce a simple yet effective \emph{Visual Probing Mechanism} (VPM), designed to identify these undesirable clusters while facilitating the recovery of semantic meaning for unequivocal cases. Using $k$-means clustering, clusters are represented by their centroids $\mu^{[C]}$,  minimizing the within-cluster cosine distance. The degree to which empirical samples resemble their assigned centroids can be quantified, in the representation space, by the cosine distance $d(x_i, \mu_{c_i})$, where $c_i = \arg \min_{c \in [C]} d(x_i, \mu_c)$ denotes the cluster assignment of $x_i$.  

        This assignment distance can be used to rank empirical samples within each cluster, which we leverage to formulate the following hypotheses:  
        \begin{itemize}
            \item \textbf{H1:} Cluster semantics can be recovered by inspecting a large number of independent samples drawn from the nearest neighbors of their centroid.
            \item \textbf{H2:} Cluster semantic consistency can be evaluated by inspecting a large number of independent samples drawn from their most distant assigned elements.
            \item \textbf{H3:} Visual inspection of the central frame of an embedding input sequence conveys sufficient information to recover its semantics.
        \end{itemize}
        While \textbf{H1} focuses on recovering the underlying semantics of clusters and assessing their granularity, \textbf{H2} aims to probe worst-case assignments to verify semantic consistency relative to the nearest samples. In particular, in the absence of supervised labels, these distant samples can serve as a proxy to assess data-fitting performance, as discussed in \Cref{sec:chapter_5_results_symbolic_eval}. In practice, we observed that inspecting the closest and furthest neighbors for $30$ different participants appear sufficient to validate \textbf{H1} and \textbf{H2}, as discussed below. Examples of assignment distances histograms are reported in \Cref{fig:chapter_5_distance_distributions}.

        \paragraph{Analysis of nearest-neighbors to recover clusters semantics and filter out ambiguous clusters (H1).\\}

        Illustrative examples of accepted and rejected clusters are shown in \Cref{fig:chapter_5_definitive_examples} and \Cref{fig:chapter_5_ambiguous_examples}, respectively. Accepted clusters correspond to highly consistent participant actions, supporting the hypothesis that the proposed VPM can be used to reliably assess the action semantics of clusters, which ultimately map to behavior. In contrast, \Cref{fig:chapter_5_ambiguous_examples} shows several ambiguous clusters that were rejected during the annotation process and are intended for refinement in later stages of the procedure.

        \begin{figure}[!htbp]
            \centering
            \includegraphics[width = 1.0\linewidth]{chapters/chapter 5/figures/chapter_5_definitive_examples.png}    
            \caption{\textbf{Illustrative examples of nearest-neighbors visual probes for nine accepted clusters.} Each subplot shows 30 independent occurrences of a given cluster, sampled from its most strongly associated nearest samples. These examples exhibit a high level of visual consistency across participants. The actions underlying each cluster can be identified unequivocally and robustly from their visual signatures. For example, (a) encodes "adding flour to the dough" and mostly occur towards the end of the assessment; (b) encodes "pulling a paper" and seem temporally uniformly distributed; (c) "opening the oven"; (d) "talking to the clinician"; (e) "mixing the dough with a "; (f) "looking at the sink"; (g) "using the mixer"; (h) "reading the recipe"; and (i) "looking at the material". We add for each cluster their (a) average assignment distance; (b) sum-of-squared cosine distance (SSD) normalized by the total SSD; (c, d) prevalence across subjects and embeddings; (e) average (s.t.d) number of segments occurrences across participants; (f) average (s.t.d.) segments duration across participants; and (g) the distribution of normalized timestamps of the drawn samples across ten equally-sized temporal bins. We note that these correspond to post-training illustrations, the statistics during training encompass training sequences only.}
            \label{fig:chapter_5_definitive_examples}
        \end{figure}

        \begin{figure}[!htbp]
            \centering
            \includegraphics[width = 1.0\linewidth]{chapters/chapter 5/figures/chapter_5_ambiguous_examples.png}    
            \caption{\textbf{Illustrative examples of nearest-neighbor visual probes for nine rejected clusters.} Each subplot displays 30 independent occurrences of a given cluster, sampled from its most strongly associated nearest neighbors. These examples illustrate ambiguous situations that, while visually correlated, exhibit varying action semantics that cannot be inferred unequivocally. Cluster-level statistics are described in the caption of \Cref{fig:chapter_5_definitive_examples}.}
            \label{fig:chapter_5_ambiguous_examples}
        \end{figure}
        

        \paragraph{Annotation image and assessment of clusters semantic consistency across nearest and furthest neighbors (H2).\\}

        We report in \Cref{fig:chapter_5_annotation_signal}-a the final annotation image for three centroids, composed of the closest and furthest neighbors visual probes, with informative statistics displayed to enrich the annotation signal. While furthest neighbors exhibit weaker level of semantic alignment with those of the nearest regions, we argue that they nonetheless provide a useful signal during evaluation to assess prototyping errors qualitatively, as discussed in \Cref{sec:chapter_5_symbolic_representation}. Finally, we observe in \Cref{fig:chapter_5_annotation_signal}-b that furthest assigned samples at the end of the recursive procedure exhibit larger homogeneity with stronger semantic alignment with the nearest neighbors, suggesting that \textbf{H2} can effectively be used to assess cluster-level semantic consistency.   % Tighten this or change to reallyevaluati...

        \begin{figure}[!htbp]
            \centering
            \includegraphics[width = 0.9\linewidth]{chapters/chapter 5/figures/chapter_5_annotation_signal.png}    
            \caption{\textbf{(a) Three illustrative examples of centroids annotation images, and (b) furthest assigned samples after model training.} Each annotation frame shows the central frame of $30$ independent nearest (left) and furthest (middle, right) neighbors from the centroids. Cluster-level statistics are described in the caption of \Cref{fig:chapter_5_definitive_examples}.}
            \label{fig:chapter_5_annotation_signal}
        \end{figure}

        \paragraph{Empirical analysis of the central-frame hypothesis (H3).\\}

        We evaluate the hypothesis that the central of the 16 frames composing the encoder input sequence can be used as a relevant visual probe to recover an embedding semantics. To this end, we observed the sequence of frames within the encoder temporal horizon for $200$ randomly sampled embeddings, and report illustrative samples of the different scenario found in \Cref{apd:chapter_5_visual_probing} with detailed discussions. Our observations suggest that the central frame is generally sufficient to represent the full sequence, except in scenarios involving extreme head or hand motion (e.g., the last rows of \Cref{fig:chapter_5_visual_probing}), where the hypothesis no longer holds. Nevertheless, we hypothesize that appearance cues and motion-blur patterns of a large number of central frames mitigate this limitation (see \Cref{fig:chapter_5_ambiguous_examples}-g for example). 

        \paragraph{Annotation of task-related versus exogeneous actions.\\}
        
        
        At each round of the iterative procedure and for each candidate cluster $c \in [C]$, the annotation signal consists in assigning one of four categories: \texttt{task-definitive}, \texttt{exogenous-definitive}, \texttt{task-ambiguous}, or \texttt{exo-ambiguous}. We notice that some actions performed during the assessment cannot readily be related to the cooking task, such as looking at the window (e.g., \Cref{fig:chapter_5_visual_probing}-b), looking down (e.g., \Cref{fig:chapter_5_visual_probing}-h), or talking to the clinician (e.g., \Cref{fig:chapter_5_definitive_examples}-d). We hypothesized that this distinction may be valuable for analyzing the behavior of patients with dysexecutive syndromes, and therefore leverage the annotation step to include the classification of the retained centroids into two disjoint categories of actions: 
        \begin{itemize}
            \item The \texttt{task-related} category, consisting of actions that are necessarily occurring for cooking the chocolate cake.
            \item The \texttt{exogeneous} category, consisting of actions that cannot readily be mapped to the cooking task, typically exhibiting non-specific temporal priors. 
        \end{itemize}

        Although some clusters fall into the ambiguous categories, we observed that the distinction between task-related and exogenous actions can still be made with high confidence in most cases. We compare the distributions of task-related and exogenous actions across diagnosis groups in \Cref{sec:chapter_6_task_exo_discrimination}.


    \subsection{Vocabulary consolidation via agglomerative hierarchical clustering}
        \label{sec:chapter_5_hierarchical_multiview_clustering}

        Since the iterative accumulation of centroids may result in semantic redundancy, where multiple centroids can correspond to the same underlying semantic concept with variations in their presentation (e.g., due to background changes, variations in ingredient packaging, or differing recipe presentations), or reflecting progressive patterns of actions leading to slightly shifted views. While such redundancy is beneficial for capturing a larger diversity of empirical samples and enhancing the robustness of our prototyping approach, it also compromises the sparsity and mutual exclusivity of the vocabulary. In particular, it inflates the vocabulary size $K$ and thus the dimensionality of subsequent symbolic representations. To address this, we aim to consolidate the vocabulary by grouping semantically similar centroids into a smaller set of representative prototypes.

        To combine the accumulated centroids into higher-level prototypes and reduce the vocabulary size, we apply a \emph{hierarchical agglomerative clustering} (HAC) to the set of accumulated clusters. HAC is a long-standing and widely used clustering method \cite{Murtagh2017, Muller2011, Mllner2013, Murtagh2011, Mllner2011ModernHA}, with successful applications in domains ranging from social network and image analysis to bioinformatics and large-scale graph mining \cite{almahamid2022agglomerative, Tie2018,  NIPS2013_35cf8659, teraHAC, Hua2017, Park2011, Tie2018, 10.1145/3097983.3098079, 10.1145/584792.584877}. Unlike the upstream clustering of \Cref{sec:chapter_5_recursive_procedure}, whose purpose was to summarize the full dataset with cluster representatives, this recombination stage focus on modeling the centroids underlying structure to facilitate gathering centroids into semantically distinct classes.

        
        While we show in \Cref{sec:chapter_5_results_consolidation} that HAC outperforms alternatives such as K-Means, DBSCAN and its hierarchical version HDBSCAN, Gaussian mixture models (GMM), and spectral clustering, it also offers interpretable representation of vision syllables as structured communities, appropriately leveraging the hierarchical structure in images. In addition, HAC can represent clusters of multiple granularity \cite{Taxonomy2014} and discover nested structures \cite{GAVRYUSHKIN2016197}. We define more formally our multi-view HAC approach in details below.

        Given a set of points $X = \{x_i\}_{i \in [n]}$, we note the set of all partitions of a set $X$ as $\mathcal{P}(X)$. A hierarchical clustering is a recursive partitioning of $X$ into a tree-structured family of nested partitions.

        The algorithm takes as input a set of $n$ points, each initially forming a singleton cluster, and proceeds iteratively by merging the most similar pair of clusters until all points are grouped into a single cluster. The procedure requires to define a dissimilarity measure between points, called $d_\mathcal{H}: \mathcal{R^D} \times \mathcal{R^D} \to \mathcal{R}^+$, as well as a dissimilarity measure between sets of points, $d_{\mathcal{L}}: \mathcal{P}(X) \times \mathcal{P}(X) \to \mathcal{R}^+$, referred to as a \emph{linkage function}. Common linkage functions are defined in terms of pairwise dissimilarities between points:  
                
        \begin{itemize}
            \item \textbf{Single linkage:} This criterion defines the distance between two clusters as the minimum dissimilarity across all cross-cluster pairs.  

            \[
                d_{\text{single}}(A,B) = \min_{x \in A, y \in B} d_\mathcal{H}(x,y).
            \]
            % While computationally simple, it may suffer from the "chaining effect", whereby clusters can be linked through narrow bridges of points, which can result in elongated or irregular cluster shapes.  

            \item \textbf{Complete linkage:} In contrast, complete linkage measures the distance between two clusters by the maximum pairwise dissimilarity, thereby encouraging the formation of compact and spherical groups.  

            \[
                d_{\text{complete}}(A,B) = \max_{x \in A, y \in B} d_\mathcal{H}(x,y).
            \]
            % While particularly sensible to the presence of outliers since single distant point can disproportionately influence the merge, this is mitigated in our setting where centroids are already carefully validated from human coding. 

            \item \textbf{Average linkage:} Here the inter-cluster distance is defined as the mean dissimilarity across all cross-cluster pairs, producing a balance between the extremes of single and complete linkage.  
            \[
                d_{\text{average}}(A,B) = \frac{1}{|A||B|} \sum_{x \in A} \sum_{y \in B} d_\mathcal{H}(x,y).
            \]
            %By averaging over all pairwise relations, this criterion is less susceptible to chaining and outliers, and typically yields clusters of intermediate shape and cohesion.  

            % \item \textbf{Ward linkage:} 
            % \[
            %     d_{\text{ward}}(A,B) = \frac{|A||B|}{|A|+|B|}\,\|\mu_A - \mu_B\|^2,
            % \]
            % where $\mu_A$ and $\mu_B$ are the centroids of clusters $A$ and $B$.  
            % This criterion minimizes the increase in within-cluster variance and is well-suited for Euclidean data.  

        \end{itemize}

        The output of HAC is a rooted binary tree, or \emph{dendrogram}, that records the sequence of merges. The leaves correspond to the initial singleton clusters, while each internal node represents a cluster obtained by merging its two children.
        
        \paragraph{Flat clusters computations from the dendrogram and estimation of the vocabulary size \(G\).\\}

            The final prototype vocabulary is defined as a partition of the original centroids by flattening the dendogram into $G$ disjoint clusters, obtained by cutting the tree into exactly $G$ disjoint subsets. The detailed procedure is described in \Cref{apd:chapter_5_flattening}. We note the partitioning function $\Phi : [K] \mapsto [G]$, which assigns each centroid index $k \in [K]$ to a prototype label $g \in [G]$. The final prototype vocabulary is noted $\nu^{[G]}_\Phi = {\nu_g}_{g \in [G]}$, with $\nu_g = \{\mu_k\}_{k \in \Phi^{-1}(g)}$.
            
            As discussed in \Cref{apd:chapter_5_model_selection_G}, determining the "optimal" number of clusters $G$ in unsupervised setting is a challenging model selection problem \cite{stop_using_elbow}. We follow established practice in unsupervised learning \cite{deAPMarques2022, LENSSEN2024102290, Batool2021, Hennig2015} and maximize the average silhouette width (\emph{silhouette score}, SS; \cite{ROUSSEEUW198753}) over a grid of candidate values of $G$. The silhouette score satisfies key axioms of clustering quality measures \cite{NIPS2008_beed1360}, making it a principled criterion for selecting the vocabulary size in our setting. We discuss this choice in more details in \Cref{apd:chapter_5_model_selection_G}.

            For a centroid $\mu_k$ assigned to prototype $\nu_g$, the silhouette score is defined as
            \[
                s(\mu_k) = \frac{b(\mu_k) - a(\mu_k)}{\max\{a(\mu_k), b(\mu_k)\}},
            \]
            where:
            \begin{itemize}
                \item \(a(\mu_k)\) is the \emph{mean intra-cluster distance} of \(\mu_k\) to all other centroids of prototype \(\nu_g\):
                \[
                    a(\mu_k) = \frac{1}{|\nu_g| - 1} \sum_{\mu_j \in \nu_g, \mu_j \neq \mu_k} d(\mu_i, \mu_j),
                \]
                where \(d(\mu_k, \mu_j)\) is the cosine distance between centroids \(\mu_k\) and \(\mu_j\), and \(|\nu_g|\) is the number of centroids in the prototype \(\nu_g\).
                
                \item \(b(\mu_k)\) is the \emph{mean nearest-cluster distance} of centroid \(\mu_k\) to all centroid of the nearest neighboring prototype \(\nu_{g'}\) (the closest prototype not containing centroid \(\mu_k\)):
                \[
                    b(\mu_k) = \min_{\nu_g' \neq \nu_g} \left\{ \frac{1}{|\nu_g'|} \sum_{\mu_j \in \nu_g'} d(\mu_k, \mu_j) \right\}.
                \]
            \end{itemize}
            The silhouette score associated with a flat clustering $\nu^{[G]}_\Phi$ is the average of \(s(\mu_k)\) over all $K$ centroids:

            $$SS(\nu^{[G]}_\Phi) = \frac{1}{K} \sum_{k=1}^K s(\mu_k),$$
                
            with values ranging from \(-1\) (incorrect clustering) to \(1\) (perfect clustering).

            The optimal number of prototype is defined as the flat clustering induced by the dendrogram cut (equivalently, by the assignment mapping $\Phi$) that maximizes the silhouette score:

            \[
            \hat{G} \;=\; \arg \max_{g \in [K]} \; SS(\nu^{[g]}_\Phi).
            \]

            This procedure requires no external supervision and, in our experiments (\Cref{sec:chapter_5_results_consolidation}), consistently produced well-defined optima with relatively high silhouette scores (e.g., $SS \geq 0.5$), indicating a good balance between \emph{within-cluster cohesion} and \emph{between-cluster separation}, and supporting the use of silhouette score maximization to consolidate centroids into stable prototypes. We omit the hat notation in the rest of the chapter to refer to the empirical prototype estimate exclusively.

        

\section{Unsupervised temporal action segmentation via kernel-based multiple change points detection}
    \label{sec:chapter_5_unsupervised_temporal_segmentation}

    The temporal dynamics of participants egocentric vision provide crucial information to compute accurate symbolic representations of behavior. Previous sections focused on prototyping the egocentric vision space to summarize dataset-level spatial appearance and motion information content, disregarding temporal dependencies and sequential context. In this section, we adopt a time series analysis (TSA) approach, where we study each participant's egocentric vision trajectory separately. Our goal is to address the key video understanding task of unsupervised temporal action segmentation (uTAS), defined above.
    
    Let a video represented as a sequence of feature vectors $X = (x_1, \ldots, x_n) \in \mathcal{X}^n \subset (\mathbb{R}^D)^n$, uTAS aims to delineate semantically consistent segments that correspond to distinct actions in the video without any labels. The set of segmentation of $[n]$ into $M$ segments is defined as $\mathrm{T}_n^M = \{0 = t_0 < ... < t_M = n\}$ for $M \in [n+1]$, with $(t_m)_{1 \leq m \leq M-1}$ the change-points between action segments. The set of all segmentation is noted $\mathrm{T}_n$.

    We discuss recent supervised and unsupervised approaches and motivate our approach in \Cref{apd:chapter_5_utas_related_work}. We propose to address uTAS by exploiting the temporal patterns in the pairwise similarity matrices, discussed in \Cref{sec:chapter_5_unsupervised_temporal_segmentation}, via the optimization of a multiple change-point objective using reproducing kernels (KCP) \cite{cpt}. This approach offers compelling properties, including (i) a strong theoretical grounding, with established consistency \cite{garreau2018consistent} and finite-sample optimality results under reasonable assumptions, taking the form of an oracle inequality in terms of quadratic risk \cite{CELISSE2018200, cpt}; (ii) the benefits of nonparametric models (e.g., inference-ready and independent of stored parameters or training data), in particular the procedure do not require to know the number of change-points in advance; and (iii) presenting efficient optimization algorithms (e.g. PELT \cite{pelt}) with fast implementations (\texttt{ruptures} package \cite{ruptures}). These are described in more details below.

    \subsection{Kernel-based multiple change point detection}
    
    % \paragraph{Introduction of the kernel theory for the change-point problem.\\}


        The generic kernel change-point detection framework was initially developed in \cite{4301363}, where detecting changes in the distribution of complex signal is reframed as the detection of mean-shifts in the sequence of points in a Reproducing Kernel Hilbert Space (RKHS). 

        Let us consider a symmetric positive semidefinite kernel $k(., .): \mathbb{R}^D \times \mathbb{R}^D \rightarrow \mathbb{R}$, and $\mathcal{H}$ the associated RKHS \cite{Aronszajn1950TheoryOR}. The canonical function $\phi: \mathbb{R}^D \rightarrow \mathcal{H}$ that maps random variable $X \in \mathcal{X} \subset \mathbb{R}$ to its embedding in the RKHS is implicitly defined by $\phi(X) = k(X, .) \in \mathcal{H}$, resulting in the following inner-product and norm: 
        $$ \forall (X_i, X_j) \in (\mathcal{R}^D)^2, \qquad <\phi(X_i) | \phi(X_j)>_{\mathcal{H}} = k(X_i, X_j) \text{  and   } \|\phi(X_i)\|^2_{\mathcal{H}} = k(X_i, X_i)$$. 

        We note for every index $1 \leq t \leq n$ the Hilbert valued random variable $Y_t = \phi(X_t)\in \mathcal{H}$. the \textit{kernel trick} allows us to express any inner product in $\mathcal{H}$ as a function of the kernel $k$, such that for every $(i, j) \in [n]^2$:
        $$\langle Y_i, Y_j \rangle_{\mathcal{H}} \;=\; k(x_i, x_j) \;=\; \mathbf{\mathcal{K}}[i,j],$$
        where $\mathbf{K}[i, j]$ denotes the coefficient $(i, j)$ of the matrix of pairwise kernel values, also called Gram matrix:  $\mathbf{K} =[k(x_i, x_j)]_{(i, j) \in [n]^2} \in \mathbb{R}^{n \times n}$.


        \paragraph{Kernel Change-point Procedure (KCP) methodology.\\}

        Let $\mathcal{P}$ denotes a Borel probability measure defined over $\mathcal{X} \subset \mathbb{R}^D$, whose indexing using $\mathcal{P}_i$ reflects the time-varying nature of the egocentric vision trajectory. The seminal kernel-based change-point were introduced in \cite{cpt}, and assume the following additive noise modeling for the latent dynamics of $Y = \{Y_i\}_{i \in [n]} \in \mathcal{H}^n$:

        \begin{equation}
            \forall 1 \leq i \leq n, \qquad Y_i = \phi(x_i) = \mu_i^\star + \epsilon_i \in \mathcal{H},
            \label{eq:latent_modeling}
        \end{equation}

        where $\mu_i^\star$ denotes the \textit{mean element} associated with the distribution $P_{x_i}$ of $x_i$, and $(\epsilon_i)_{i \in [n]}$ independent and identically distributed (i.i.d), with 0 mean and variance $\sigma^2$.

        If $\mathcal{X}$ is separable and the kernel verify $\mathbb{E}_{X \sim \mathcal{P}}[k(X,X)] < +\infty$ \cite{Ledoux1991}, which holds for $\mathcal{X} \subset \mathcal{R}^D$ and the typical kernels employed in this chapter (e.g. linear and Gaussian kernels), then this mean element exists and is unique, defined as: $\mu^\star_i = \mathbb{E}_{X \sim \mathbb{P}_{i}}[\phi(X)]$.   

        If the chosen kernel is characteristic, meaning that the mean embedding map $\mathcal{P} \mapsto \mu^\star$ is injective \cite{Sriperumbudur}, then distinct distributions yield distinct mean elements. Formally,
        \begin{equation}
            \forall\, 1 \leq i \neq j \leq n, 
            \qquad \mathcal{P}_{i} \neq \mathcal{P}_{j} 
            \;\; \implies \;\; \mu^\star_i \neq \mu^\star_j.
            \label{eq:injective}
        \end{equation}

        The contraposition of \Cref{eq:injective} underlies the central idea of kernel change-point detection: action boundaries can be detected as mean shifts in the latent RKHS trajectory $(Y_i)_{i \in [n]}$ \cite{cpt}. Said differently, this model assumes that the egocentric vision embeddings are mapped via $\phi$ into a signal in $\mathcal{H}$ with piecewise-constant mean, i.e., that the distributions $\big(\mathcal{P}_i \big)_{i \in [n]}$ are stationary within each segment.


        The quadratic risk associated with a segmentation $\tau^{M} = (t_m)_{1 \leq m \leq M} \in \mathrm{T}_n^M$ is defined as \cite{4301363}:
        \begin{equation}
            \hat{\mathcal{R}}_n(\tau^{M}) 
            = \tfrac{1}{n}\sum_{m=0}^M \sum_{i=t_m}^{t_{m+1}-1} 
            \big\| \phi(x_i) - \bar{\mu}_{t_m:t_{m+1}} \big\|_{\mathcal{H}}^2,
            \label{eq:opti}
        \end{equation}
        where $\bar{\mu}_{t_m:t_{m+1}}$ denotes the empirical mean of the segment $\{\phi(x_{t_m}), \ldots, \phi(x_{t_{m+1}-1})\}$ (since $\mu^\star$ is unknown). 
        
        %This expression can be rewritten using the kernel trick as:

        % \begin{equation}
        %     \hat{\mathcal{R}}_n(\tau) = \sum_{m=0}^M \Bigg[ 
        %             \sum_{i=\tau_m}^{\tau_{m+1}-1} \mathcal{K}[i,i]
        %             - \frac{1}{\tau_{m+1}-\tau_m} 
        %             \sum_{i,j=\tau_m}^{\tau_{m+1}-1} \mathcal{K}[i,j]  \Bigg].
        %     \label{eq:opti-kernel}
        % \end{equation}

        \paragraph{Penalization term for unknown number of change-points.\\}
        
        Since the number of change points is unknown, a regularization term is typically considered to penalize model complexity $\text{pen}: \mathrm{T}_n \to \mathbb{R}$. The choice of a suitable penalization is a classical but difficult model selection problem \cite{Lavielle2005,Brodsky1993}. Existing approaches include linear or nonlinear constraints on the number or positions of change points \cite{Verzelen2023, cpt, Zhang2007}, popular information criteria (AIC \cite{Akaike1998}, Mallows’ $C_p$ \cite{Mallows1973}, BIC \cite{Schwarz1978}), $k$-fold cross-validation \cite{Golub1979}, and stability analysis \cite{Liu2010, Meinshausen2010}. A more extensive discussion with additional references is provided in \cite{cpt}. 

        In our experiments, we adopt a penalty shape proportional to the number of change point (similar to AIC, BIC and $C_p$ criteria), such as the penalized objective becomes \cite{garreau2018consistent}:

        \begin{equation}
            \hat{\mathcal{R}}_{n, \lambda}(\tau^M) 
            = \tfrac{1}{n} \sum_{m=0}^M \sum_{i=t_m}^{t_{m+1}-1} 
            \big\| \phi(x_i) - \bar{\mu}_{t_m:t_{m+1}} \big\|_{\mathcal{H}}^2 + \tfrac{\lambda}{n} |t|, 
            \label{eq:opti-pen}
        \end{equation}

        with $\lambda > 0$, controlling the trade-off between data fidelity and number of change points. This penalty is proven to be optimal in the sense of the \textit{unbiased risk estimation heuristics} for $\lambda = 2 \sigma^2$ (where $\sigma$ is the variance of the additive noise $\epsilon_i$ in \Cref{eq:latent_modeling}; see more information in section 2.2 in the recent survey \cite{slop_heuristics_survey}). 
        
        In order to estimate this optimal multiplicative constant $\lambda$, we adopt the \emph{minimum penalty} principle \cite{slop_heuristics_survey}, specifically the \emph{slope heuristic}, which is well-established in the model selection literature \cite{Baudry2011, Lavielle2005, JMLR:v10:arlot09a, lebar2002} and have been applied to kernel change-point detection \cite{CELISSE2018200, cpt, NIPS2008_08b255a5}. We use the procedure of the recent survey \cite{slop_heuristics_survey} (see alg. 2 in Section 2.4). While the vanilla kernel change point detection (KCP) framework determine an optimal penalty term whose coefficients is estimated for each time series, this is undesirable in our case where we aim to perform fair comparisons across the participant's symbolic representations. We thus adapt it using a fixed-effect model to estimate the slope using a participant-level intercept \cite{10.5555/1162264}. This allows to estimate a single penalty common to all participants, ensuring a fair comparison between participants. 

        
        Specifically, we first estimate the slope $\hat{\beta}$ of $(\mathcal{R}_{N_i}({\tau^{m}})_{i\in[1, N_{\text{train}}]}$ as a function of $m$ in the typical range $m \in [0.6 M_{\text{max}}, M_{\text{max}}]$, using ordinary least squares regression with group-level intercepts (fixed-effect design). $M_{\text{max}}$ is chosen "large enough", that is, $M_{\text{max}} = 900$, which corresponds to a change point every 2 seconds of the median video duration ($\approx 30$ minutes). The optimal penalty for each participant is then defined as $\hat{\lambda}_{\text{opt}} = -2n*\hat{s}$. In practice, we observed that using about $30$ values of $m$ within the typical range $[0.6 M_{\text{max}}, M_{\text{max}}]$ were sufficient for robustly estimating the slope ($R^2 =.99$). The detailed optimization procedure to solve \Cref{eq:opti} using dynamic programming can be found in \cite{cpt}, and is implemented in the \texttt{ruptures} package \cite{ruptures}. %CODE: TODO:, as illustrated in \Cref{fig:chapter_5_slope_heuristic}.

        Finally, for each participant's sequence $i$, we compute the optimal temporal segmentation, $\hat{\tau}_{i}$, by solving the penalized optimization problem of \Cref{eq:opti-pen} with $\lambda = \hat{\lambda}_{\text{opt}}$. The optimization is performed using the Pruned Exact Linear Time (PELT) algorithm \cite{pelt}, which leverages an optimality-based pruning rule to discard suboptimal change-point candidates during the dynamic programming recursion. While its worst-case computational complexity remains quadratic, $\mathcal{O}(n^2)$, the pruning rule ensures an expected linear runtime, $\mathcal{O}(n)$ under mild regularity conditions, and is efficiently implemented in the \texttt{ruptures} library \cite{ruptures}.


\section{Computation of the symbolic representation and evaluations}
    \label{sec:chapter_5_symbolic_representation}


    The last component of our symbolic representation method combines the estimated prototypes vocabulary, $\nu^{[G]}_\Phi \in \mathcal{P}(\mu^{[K]})$, with the action boundary detections for each participant's sequence, $\tau_{i} \in \mathrm{T}_{N_i}$, for $i\in [N_{subjects}]$.


    \paragraph{Prototypes inference and filtering.\\}


    Given the latent embedding sequence of a participant, $\textbf{x} = (x_i)_{i \in [n]} \in (\mathbb{R}^D)^n$, we start by assigning each sample its closest centroid index, $\mathbf{c} = (c_i)_{i \in [n]} \in [K]^n$, with $c_i= \arg \min_{k \in [K]}  d(x_i, \mu_k)$, and note their assignment distance $\mathbf{d} = (d_i)_{i \in [n]}$ with $d_i =  \min_{k \in [K]}  d(x_i, \mu_k)$. The reduction mapping $\Phi$ and centroids assignments are then use to compute prototypes assignments, $\mathbf{y} = \{ y_i \}$ with $y_i = \Phi(c_i)$.
    
    In order to improve symbolic representations accuracy, we withdraw the assignment of furthest samples to their prototypes if their assignment distance is over a centroid-specific threshold. In practice, we observed that using one standard deviation over the centroids' average assignment distance yielded visual probes of furthest assigned sample (cf $\mathbf{H2}$) more homogeneous and consistent with the class semantic (this evaluation is reported in \Cref{sec:results_generalizability_fidelity}). We leave for future work the development of more tailored identification of sample-prototypes mismatch. Specifically, we note $\textbf{d}_{\nu_j}$ the set of assignment distances of a prototype $j\in [G]$, defined as:

    $$\textbf{D}_{g} = \{ \min_{\mu_k \in \nu_g} \big( d(x, \mu_k)\big) \}_{x \in V_g},$$
    
    where $V_g = \{x_i \;\mid\; y_i = g \}$ defines the influence area of prototype $g$, also called Voronoi cell. The mean and standard deviation of prototypes assignment distances are noted $\textbf{m} = (m_g)_{g \in [G]}$ and $\mathbf{\sigma} = (\sigma_g)_{g \in [G]}$. %, computed empirically as $\hat{m}_{\nu_j} = \tfrac{1}{| \mathcal{V}_j| }\sum_{x \in \mathcal{V}_j} \min_{k \in \Phi^{-1}(j)} \big( d(x, \mu_k)\big)$.

    The filtered prototypes assignments are then computed as $\tilde{y} = (\tilde{y}_i)_{i \in [n]} \in ({-1}\cup[G])^n$ with $\tilde{y_i} = -\mathbbm{1}(d_i > m_{\nu_{y_i}}+ \sigma_{\nu_{y_i}} ) + \mathbbm{1}(d_i \leq m_{\nu_{y_i}} +  \sigma_{\nu_{y_i}} ) y_i$, such that the label \texttt{-1} encodes a "background" class that reflects assignment uncertainties.
    

    \paragraph{Segments label computation.\\}

        We finally use the temporal segmentation predictions to assign a unique label per segment by applying a majority voting criteria restricted to the two central quartiles of each segment only. Indeed, we hypothesized that the samples near boundaries may predict prototype assignment less reliably, due to motion patterns or visually unstable transition phases causing poor encoding and spurious association with prototypes. This was verified empirically in preliminary experiments where we show that segments' border quartiles exhibit significantly higher assignment distances compared to the central quartiles.
        
        Segment-level assignment improves the robustness of our symbolization procedure by smoothing the action sequence to reflect a more realistic visual temporal regularity, as illustrated in \Cref{sec:chapter_5_results_temporal_segmentation}. 

        Formally, let a video embedding sequence be denoted as  $x = \{x_i\}_{i \in [n]}$, with estimated temporal segmentation $\tau = (t_0, \ldots, t_M) \in \mathrm{T}_n$, and pointwise prototype assignment labels $\tilde{\mathbf{y}} = \{\tilde{y}_i\}_{i \in [n]}$. The piecewise constant symbolic representation is obtained through the following steps:

        \begin{enumerate}
            \item \textbf{Segment-level aggregation.}
            For each segment $j \in [M]$, defined by boundaries $[t_{j-1}, t_j]$, let its length be $\ell_j = t_j - t_{j-1}$.
            We define its the central interval as
            \[
                Q_j = 
                \big[
                    t_{j-1} + \tfrac{1}{4}\ell_j,\;
                    t_j - \tfrac{1}{4}\ell_j
                \big]
                \subset [t_{j-1}, t_j],
            \]
            corresponding to the two central quartiles of the segment.

            \item \textbf{Label assignment.}
            The label of segment $j$ is computed by majority voting within 
            its central interval:
            \[
                \bar{s}_j = 
                \operatorname{MajVote}\big(
                    \{\tilde{y}_i : i \in Q_j\}
                \big),
            \]
            where $\operatorname{MajVote}$ denotes the majority voting operator.

            \item \textbf{Piecewise constant reconstruction.}
            The symbolic representation of the sequence is defined as
            \[
                \mathbf{s} = \{s_i\}_{i \in [n]},
                \quad 
                s_i = \bar{s}_j
                \quad \text{for } 
                t_{j-1} \le i < t_j.
            \]

            \item \textbf{Temporal smoothing.}
            Finally, the symbolic representation is refined by removing any short segments that are surrounded by the same label on both sides, provided the segment length is less than 4 embeddings ($\approx 1.3s$).
        \end{enumerate}

        Finally, we note that this symbolization procedure allow to prune the temporal segmentation predictions by removing change-points between identical segment labels, which is studied in the result section.
        




\section{Experimental setting}
\label{sec:experimental_setting}

    \subsection{Model estimation procedure}

        In order to assess the generalizability of our prototyping approach, we divide the dataset into a training set, $D_{\text{train}} = \{ \textbf{x}_s  \}_{s \in \mathcal{I}_{\text{train}}}$, and a testing set, $D_{\text{test}} = \{ \textbf{x}_s  \}_{s \in \mathcal{I}_{\text{test}}}$, with $\mathcal{I}_{\text{train}}$ (respectively, $\mathcal{I}_{\text{test}}$) containing $80\%$ (respectively, $20\%$) of randomly selected participants. That is, the latent sequence of 129 participants ($N_{\text{train}}=725335$ embeddings) were used to estimate prototypes and calibrate the penalty parameters, $\hat{\lambda}_{\text{opt}}$, and 33 participants sequences ($N=180568$ embeddings) are used to evaluate the symbolic representation performances. 

        We collected $K_\text{in}=800$ centroids across $P=8$ rounds during which $C=100$ centroids were estimated at each round followed by the feedbacks collections stage. Our stopping criteria is based on the following heuristics: we aimed to collect a significantly higher number of centroids than the average number of action segments detected during the unsupervised temporal segmentation step, that is $\bar{M}=274$ $(STD=24)$ change points. We aimed more for a proof of concept here, noting that a larger number of rounds could benefit model exhaustivity and accuracy, with limited impact on final prototypes cardinality due to the consolidation step. 

    \subsection{Implementation details}

        Cosine $k$-means clustering were performed using the \texttt{Faiss} library \cite{douze2024faiss} with default parameters. Annotation were performed using the lightweight \texttt{pigeon-annotation} software\footnotemark[1]. The NMI and ARI  $95\%$ confidence interval (CI) were estimated via 1000 bootstrap resamples. %Optimal parameters of the inverse sigmoid function $(\hat{s},\hat{\lambda_0})$ are estimated via non-linear least squares, as implemented by the \texttt{curve_fit} function of \texttt{scipy} Python package (v1.10.1).
        \footnotetext[1]{\url{https://github.com/agermanidis/pigeon/tree/master}}

        

\section{Evaluation protocol and results}
\label{sec:chapter_5_evaluation_protocol}

    \subsection{Evaluation of action boundaries detection granularity}
    \label{sec:chapter_5_results_temporal_segmentation}

        We detect transitions between piecewise constant regimes in the embedding dynamics using a kernel-based multiple change-points detection algorithm. While the exact optimization guarantees the recovery of all transition between hypothetically constant regimes, the granularity, or magnitude of the "mean-shift" between regimes, is controlled by a global regularization parameter, $\hat{\lambda}_{\text{opt}}$, which penalizes the introduction of change-points at the same level across participants' sequences.
        
        In the absence of fine-grained temporal annotations, which precludes quantitative evaluation with supervised metrics, we exploit participants’ Gram matrices as a rich visualization of the whole session dynamics. To assess whether the detected change-points reflect the desired finest action granularity, we examine the superposition of change-point results with Gram matrices, verifying that the segmentation captures all smallest diagonal blocks. Finally, we analyze the empirical distribution of segment durations to confirm their alignment with naturalistic vision segments.

        \subsubsection{Results}
        %Our aim is to estimate an ethologically relevant segmentation of participants egocentric videos into fine-grained semantically consistent motifs, with a temporal quantification that is comparable across participants.
        Visual examples of estimated actions segmentation for six participants are shown in \Cref{fig:gram_participant_cpts}, suggesting that the detected change-points accurately capture fine-grained transitions corresponding to the smallest diagonal blocks in the Gram matrices. In addition, the estimated temporal segments exhibited a realistic large range of durations, between 1.6 and 21.4 seconds at the $5^{th}$ and $95^{th}$ percentiles. 

        \begin{figure}[!htbp]
            \centering
            \includegraphics[width=1.0\linewidth]{chapters/chapter 5/figures/gram_participant_cpts.png}    
            \caption{\textbf{Visual examples of six participants' Gram matrices with detected change-points overlaid.} Detected change-points (vertical and horizontal lines) are shown on each Gram matrix, suggesting that the segmentation accurately captures fine-grained actions transitions, isolating small segments of similar pairwise cosine similarity. For example, the initial diagonal block of participant \#5 corresponds to the patient reading successive pages of the recipe, which is consistently appearing during the rest of the session, and correctly detected as such by the symbolic representation. Within this block, three regions exhibit characteristic harlequin patterns (green boxes), caused by the participant’s hand intermittently entering the egocentric view. These local fluctuations were not isolated as distinct segments, due to the optimization procedure evaluating their within-segment variability as sufficiently low to estimate a single segment, illustrating the role of the \textit{minimum penalty} heuristics in selecting a suitable penalty parameter to prevent over-segmentation.}
            \label{fig:gram_participant_cpts}
        \end{figure}

        This confirms the reliability of the temporal segmentation component, which underpins the segment-level majority-voting assignment of prototypes. By enforcing consistency within segments, this step enhances both the robustness and temporal smoothness of the symbolic representation pipeline. An illustrative example is provided in \Cref{fig:chapter_5_temporal_smoothing}, showing how segment-level assignment reduces spurious label fluctuations and yields more regular symbolic trajectories.

        \begin{figure}[!htbp]
            \centering
            \includegraphics[width=0.9\linewidth]{chapters/chapter 5/figures/chapter_5_temporal_smoothing.png}    
            \caption{\textbf{Effect of temporal smoothing on symbolic representations.} 
            Symbolic sequences from 10 randomly sampled egocentric videos are shown using raw prototype assignments (a) and after applying unsupervised temporal action segmentation with segment-level majority voting (b). The smoothing step reduces noise and yields more temporally coherent prototype sequences.}
            \label{fig:chapter_5_temporal_smoothing}
        \end{figure}

    \subsection{Evaluation of the centroids-to-prototypes consolidation stage}
    \label{sec:chapter_5_results_consolidation}

        The egocentric vision prototypes are derived as the leaf subtrees of a hierarchical agglomerative clustering (HAC) applied to the centroids accumulated during the recursive latent-space exploration procedure (\Cref{sec:chapter_5_recursive_procedure}). While individual centroids are manually filtered to ensure they provide unequivocal encodings of specific actions, their aggregation through HAC may occasionally merge centroids corresponding to distinct actions, thereby yielding semantically ambiguous prototypes. To evaluate the semantic consistency of the consolidated prototypes, we employ a visual probing strategy that is similar to the one used during centroid-level filtering. Specifically, each prototype is manually inspected and assigned to either \texttt{valid} or \texttt{ambiguous} categories according to stringent consistency criteria. For each prototype, a visual probe is generated by concatenating the visual probes of its constituent centroids, as illustrated in \Cref{fig:chapter_5_prototype_visual_probe}. Since the primary objective of the consolidation step is to reduce redundancy among centroids, a prototype is rejected if any pair of its constituent centroids is judged to represent distinct underlying actions.

        \begin{figure}[ht]
            \centering
            \includegraphics[width=1.0\linewidth]{chapters/chapter 5/figures/prototype_visual_probe.png}
            \caption{\textbf{Illustration of a prototype visual probe.} Each item corresponds to the nearest-neighbors visual probe of a centroid (see \Cref{sec:chapter_5_annotation}). In this example, the prototype is accepted, as all constituent centroids consistently capture the same underlying action of "pouring an ingredient", noting that the distinctions between the type of ingredient (e.g. sugar, floor) is not currently captured by our approach.}
            \label{fig:chapter_5_prototype_visual_probe}
        \end{figure}


        \subsubsection{Results}

        After the $p=8$ rounds of centroid accumulation, annotators retained $K_1 = 642/800$ and $K_2 = 613/800$ centroids for \texttt{Model-A} and \texttt{Model-B}, respectively. To investigate the temporal emergence of final prototypes during the iterative accumulation procedure, we analyzed the provenance of centroids contributing to each prototype across the eight rounds. For each vocabulary, we show in \Cref{fig:chapter_5_emergence_prototypes} the cumulative earliest round of appearance of the prototypes constituents. While it shows that the majority of prototypes appear during the first iteration, novel prototypes are effectively being discovered as the procedure progresses. Prototypes incorporated centroids from an average of $3.58 \pm 2.14$ rounds for \texttt{Model-A} and $3.39 \pm 2.10$ rounds for \texttt{Model-B}, demonstrating the centroids accumulation redundancy across rounds.

        \begin{figure}[!htbp]
            \centering
            \includegraphics[width = 0.9\linewidth]{chapters/chapter 5/figures/chapter_5_emergence_prototypes.png}    
            \caption{\textbf{Cumulative earliest round of appearance of prototypes constituents for each model.} While the majority of prototypes emerge during the first iteration, novel prototypes are effectively discovered as the procedure progresses, suggesting the iterative refinement's capacity to uncover novel concepts.}
            \label{fig:chapter_5_emergence_prototypes}
        \end{figure}
        
        The prototype vocabulary was then defined by flattening the hierarchical clustering at the level maximizing the silhouette score across increasing numbers of clusters, thus achieving the best compromise between within-prototype consistency and across-prototype separability. As shown in \Cref{fig:chapter_5_silhouette_score}-b, the silhouette score exhibited a similar profile and a well-defined optimum for both models, yielding $G = 178$ prototypes each.

        \begin{figure}[!htbp]
            \centering
            \includegraphics[width=0.9\linewidth]{chapters/chapter 5/figures/chapter_5_silhouette_score.png}    
            \caption{\textbf{Evolution of silhouette score across number of estimated prototypes.} The optimal number of prototypes is identified as the one maximizing the silhouette score, yielding $G=178$ prototypes for \texttt{Model-A} and \texttt{Model-B} during the $G=75$ prototypes for \texttt{Model-A} and $G=65$ for \texttt{Model-B}. We note that the predicted number of clusters for DBSCAN and H-DBSCAN depends on the neighborhood size and minimum number of samples per clusters fixed in the experiments ($\epsilon=0.5$ and $\text{min}\_\text{sample}=5$).}
            \label{fig:chapter_5_silhouette_score}
        \end{figure}


        To address the redundancy that remained among these prototypes, we reapplied the HAC algorithm on the mean-aggregated centroids, reducing the vocabulary cardinality to $G = 75$ for \texttt{Model-A} and $G = 65$ for \texttt{Model-B} (see \Cref{fig:chapter_5_silhouette_score}-c). Manual inspection of the resulting prototypes confirmed that the majority were valid (see e.g. \Cref{fig:chapter_5_visual_probes}), while rejected prototypes mainly corresponded to merges of centroids encoding distinct actions that share similar visual features. We report examples of valid and rejected prototypes in \Cref{fig:chapter_5_prototype_visual_probe} and \Cref{apd:chapter_5_prototype_visual_probe_additional}, and discuss consequences and possible improvements in \Cref{sec:chapter_5_conclusion}. The full progression of prototype estimation,  acceptance rates, as well as the unsupervised symbolic representation measures are summarized in \Cref{tab:chapter_5_prototype_pipeline}. We visualize in \Cref{fig:chapter_5_final_symbolic_representation} the final symbolic representations obtained for all participants, and provide additional qualitative analyses in \Cref{apd:chapter_5_prototype_additional}.


        \begin{table}[!htbp]
            \centering
            \begin{tabular}{lcc}
                \toprule
                Stage & \texttt{Model-A} & \texttt{Model-B} \\ 
                \midrule
                Initial centroids ($p=8, C=100$) & 800 & 800 \\
                Valid centroids  & 642 & 613 \\
                $1^{st}$ HAC & 178 & 178 \\
                $2^{nd}$ HAC & 75 & 65 \\
                \midrule
                Accepted prototypes & 55 (72.3\%) & 45 (69.2\%) \\
                Rejected prototypes & 21 (27.7\%) & 20 (30.8\%) \\

                \midrule
                Background prop. (\%) & 15.7 $\pm$ 7.6 & 15.8 $\pm$ 7.7 \\
                Segment duration (s) & 8.357 $\pm$ 12.437 & 8.449 $\pm$ 12.878 \\
                $p^{\text{labels}}_{\text{change}}$ $(\%)$ & 29.5 $\pm$ 3.8 & 28.0 $\pm$ 4.0 \\
                $p^{\text{withdrawn}}_{\text{CP}}$ $(\%)$ & 21.5 $\pm$ 4.7 & 22.4 $\pm$ 5.2 \\
                $d_\text{tot}$ & 1856.5 $\pm$ 622.9 & 1840.3 $\pm$ 608.5 \\
                \bottomrule
            \end{tabular}
            \caption{Prototype vocabulary estimation pipeline for \texttt{Model-A} and \texttt{Model-B}, from initial centroid accumulation to final acceptance after redundancy reduction and manual inspection. $p^{\text{labels}}_{\text{change}}$: proportion of label switches between pointwise assignments and their segment-level consensus, $p^{\text{withdrawn}}_{\text{CP}}$: proportion of change-points occurring between two identical consecutive labels, $d_\text{tot}$: total sum of cosine assignment distance.}
            \label{tab:chapter_5_prototype_pipeline}

        \end{table}
        % CODE: TODO: add number of segments

        \begin{figure}[!htbp]
            \centering
            \includegraphics[width=0.8\linewidth]{chapters/chapter 5/figures/chapter_5_final_symbolic_representation.png}    
            \caption{\textbf{Symbolic representations of participant's egocentric vision using \texttt{Model-A} (a) and \texttt{Model-B} (b).} Each row corresponds to a different participant performing the chocolate cake cooking task, and each color represents a different cluster. Color-coding is independent for each model. Sequences are upsampled to match the longest sequence for visualization purposes.}
            \label{fig:chapter_5_final_symbolic_representation}
        \end{figure}

        In the rest of the evaluation, we mitigate the influence of ambiguous prototypes by comparing the vocabularies and symbolic representations using \texttt{valid} prototypes only.


    \subsection{Qualitative evaluation of symbolic representations fidelity and generalizability}
        \label{sec:chapter_5_results_symbolic_eval}

        The fidelity of our symbolic representations in capturing participants’ behavior is essential for ensuring meaningful and reliable clinical interpretations. In particular, we seek to evaluate whether empirical samples are consistently and accurately aligned with their assigned prototypes. However, the absence of ground-truth labels prevents the use of supervised metrics to assess temporal precision or behavioral classification accuracy at fine-grained levels. We therefore adopt complementary qualitative and proxy quantitative evaluations, focusing on both fidelity (accuracy of prototype assignments) and generalizability (consistency across participants).  

        \begin{itemize}
            \item \textbf{Generalizability of the vocabulary.}  
            We assess whether the egocentric vision vocabulary extends to unseen participants. For each prototype, we examine the \textit{visual semantics} consistency between training and testing participants using the nearest-neighbor visual probes described in \Cref{sec:chapter_5_annotation}.  

            \item \textbf{Data fidelity.}  
            As a proxy for assignment accuracy, we leverage the empirical relationship between assignment distance and semantic consistency, and focus the assessment on weakest assignment. Specifically, we inspect furthest-neighbor visual probes to test whether samples with the largest assignment distances remain consistent with the action identity established by nearest neighbors.  

            \item \textbf{Within- and across-participant consistency.}  
            For each prototype, we visualize up to ten segments occurrences across twenty participants to assess (i) the retrieval of multiple instances of the same action \textit{within} individual sequences, and (ii) the semantic validity of prototype assignments \textit{across} different participants, in both training and test sets.  
        \end{itemize}

        For each of these experiments, we randomly sample ten visual probes, reporting five representative examples in \Cref{fig:chapter_5_visual_probes}, with five additional probes provided in \Cref{apd:chapter_5_visual_probes_second}.

        \subsubsection{Results}
        \label{sec:results_generalizability_fidelity}

        \paragraph{Generalizability of the prototypes.\\}

        The estimated prototypes were found to transfer remarkably well to egocentric videos of participants unseen during training. Importantly, the visual semantics of prototypes, as assessed through nearest-neighbor visual probes, remained stable across both training and testing participants, as illustrated in \Cref{fig:chapter_5_visual_probes}.

        \begin{figure}[!htbp]
            \centering
            \includegraphics[width=0.9\linewidth]{chapters/chapter 5/figures/chapter_5_visual_probes.png}    
            \caption{\textbf{Randomly selected visual probes of egocentric video prototypes.} Prototype generalizability to testing participants is assessed via nearest-neighbor visual probes (left and middle columns). Semantic consistency within prototypical regions is further examined in testing participants by comparing nearest- and furthest-neighbor visual probes (middle and right columns). Prototypes were randomly sampled from \texttt{Model-B}.}
            \label{fig:chapter_5_visual_probes}
        \end{figure}

        \paragraph{Fidelity assessment.\\}

        Analysis of furthest-neighbor probes revealed underlying actions that remained semantically consistent with those identified in nearest-neighbor regions. This demonstrates that regions in the videoMAE-v2 latent space stratified by centroids' assignment distance preserve coherent egocentric semantics, providing an empirical validation of the fidelity of symbolic representations. Notably, even samples with the weakest assignment to a prototype tended, in most cases, to remain aligned with the prototype’s semantic class. This supports the use of our assignment distance-based filtering heuristic: discarding samples beyond one standard deviation of the prototype assignment distance appears a reasonable trade-off between precision and pruning. Future work could explore more sophisticated identification of misrepresented samples to improve the accuracy of the symbolic representations.

        We further evaluate the prototypes accuracy by visualizing multiple occurrences of each prototype across participants (\Cref{fig:chapter_5_subjects_rows}). The consistent capture of the same underlying action both within- and across-individuals highlights the robustness of our similarity-based approaches and the feasibility of egocentric vision prototyping. Additional examples are provided in \Cref{apd:chapter_5_visual_probes_second}.

        \begin{figure}[!htbp]
            \centering
            \includegraphics[width=0.8\linewidth]{chapters/chapter 5/figures/chapter_5_subjects_rows.png}    
            \caption{\textbf{Illustrative examples of prototype assignments within and across participants.} For each prototype, up to ten segment occurrences are sampled across twenty participants. While occasional outliers are present, the majority of samples consistently capture similar actions across individuals, highlighting both the robustness and the semantic coherence of the prototype vocabulary.}
            \label{fig:chapter_5_subjects_rows}
        \end{figure}

        Despite this overall consistency, both types of visual probes reveal the presence of outliers that appear spuriously assigned to prototypical regions. These mis-assignments may arise from (i) action's transition phases yielding semantically ambiguous samples with respect to the vocabulary, (ii) the use of the central frame to recover embeddings' interpretability (see \Cref{apd:chapter_5_visual_probing}), (iii) the vocabulary's lack of exhaustivity, or (iv) the imperfection of samples encoding. We discuss their impact and potential improvements in the discussion section.  

        \paragraph{Exhaustivity of the vocabulary.\\} 

        The estimated vocabularies were found to cover a wide range of egocentric actions, with no prototypes appearing specific to either \texttt{Model-A} or \texttt{Model-B}. This suggests that the prototyping procedure effectively captures the diversity of actions present in the dataset. This is further supported by the two-dimensional projection of prototypes and empirical samples in \Cref{fig:chapter_5_umap_visualization}, where prototypes seem to cover the majority of empirical samples, suggesting an effective summarization of participants' visual experiences. 
        
        \begin{figure}[!htbp]
            \centering
            \includegraphics[width = 0.9\linewidth]{chapters/chapter 5/figures/chapter_5_umap.png}    
            \caption{\textbf{Two-dimensional UMAP projection of the prototype vocabulary centroids, training samples (top) and test samples (bottom).} The prototype vocabulary spans a substantial portion of the empirical training and testing samples, indicating broad coverage of participants' visual experience during the cooking task. This low-dimensional projection reveals the proximity between visually similar action classes (e.g., opening the fridge (red) and manipulating the oven (brown)) and temporally close actions (e.g., preparing the mixer (dark purple) and using it (light green)).}
            \label{fig:chapter_5_umap_visualization}
        \end{figure}
        
        However, certain rare actions, such as "washing hands in the restroom", were not consistently captured by the prototypes. This may be attributed to their lower frequency in the dataset or their inherent variability in execution across participants. Future work could explore strategies to enhance the representation of such actions, potentially through refined residual space definition, or using a larger training dataset. 
        
        Lastly, we note that some actions were not distinctly represented, such as "pouring sugar" and "pouring flour". This is due to their similarity in visual appearance and motion signatures, suggesting that (i) the iterative refinement procedure failed to disentangle these cases using the current residual spaces definition, and (ii) a lack of discriminative power of the video encoder. This highlights a limitation in the current approach's ability to differentiate between visually similar but semantically distinct actions.

        %We note that in a preliminary experiment, where all participants data was used for prototyping yielded similar results, at the exception that a prototype further supporting the robustness of our approach.

    \subsection{Sensitivity to the human feedback signal}
        \label{sec:chapter_5_results_human_feedback}

        To assess the robustness of our prototyping procedure to subjective variability in human feedback, we created two independent prototype models, denoted \textsc{Model-A} and \textsc{Model-B}, based on annotations from two different raters during the prototyping process described in \Cref{sec:chapter_5_prototyping}. Prior to annotation, a short coaching session aimed to explain the annotation protocol described in \Cref{sec:chapter_5_annotation}.
        
        \subsubsection{Inter-rater reliability} 

            We first quantified agreement between the two annotators on a shared set of $400$ images. Inter-rater reliability was measured using Cohen’s $\kappa$ coefficient \citep{cohen1960}, which adjusts observed agreement for chance-level concordance based on the marginal label distributions. Interpretation of $\kappa$ values follows established guidelines \cite{McHugh2012-bu, Landis1977}. To account for sampling variability, we performed a non-parametric bootstrap with $1{,}000$ resamples (with replacement) of the $400$-image set, yielding an empirical $95\%$ confidence interval for $\kappa$. To disentangle the sources of disagreement, we compute the annotators’ confusion matrix under three label grouping schemes: (\emph{i}) the original four fine-grained categories (\texttt{task-ambiguous}, \texttt{task-definitive}, \texttt{exo-definitive}, \texttt{exo-ambiguous}); (\emph{ii}) a binary grouping by cluster quality (\texttt{ambiguous} vs.~\texttt{definitive}); and (\emph{iii}) a binary grouping by cluster semantic (\texttt{task} vs.~\texttt{exogeneous}). This allowed us to determine whether reliability was primarily influenced by subjective judgments of prototype quality (certainty) or by the task/exogenous distinction.
            Together, these analyses provide both a quantitative reliability estimate and a qualitative characterization of the annotation process, thereby strengthening the validity of subsequent comparisons between \textsc{Model-A} and \textsc{Model-B}.

            \paragraph{Results.\\}

                The two annotators exhibited \textit{moderate} agreement when classifying a shared set of $400$ images, with Cohen's $\kappa$ values of $0.46$ [95\% CI: 0.37, 0.55] for the full four-category scheme, $0.44$ [95\% CI: 0.35, 0.54] when grouping by certainty (ambiguous vs.~definitive), and \textit{substantial} agreement of $0.78$ [95\% CI: 0.57, 0.93] when grouping by semantic category (task-related vs.~exogenous). These results indicate that while there is substantial concordance on the task/exogenous distinction, judgments about prototype quality seem more subjective. Discrepancies were found in clusters where experience with the procedure matters: some visual situation are doomed to be stick together using the videoMAE-v2 encoder and our recursive procedure (for example, pouring ingredients with different packaging were not disentangled; see more discussions in \Cref{sec:chapter_5_results_symbolic_eval}), leading to higher acceptance rate for the author compared to the second annotator. We believe higher level of raters agreements can be achieved in future experiments, by improving coaching session design, clarifying the exogenous category, and providing more training examples.

                Next, we study whether these differences led to significant differences in the final prototype vocabularies, as well as in the resulting symbolic representations.

                \begin{table}[ht]
                    \centering
                    \caption{Inter-annotator agreement measured with Cohen's $\kappa$ for different category grouping schemes. Values in brackets indicate the 95\% confidence interval obtained by bootstrap resampling.}
                    \begin{tabular}{l c}
                    Grouping & $\kappa$ [95\% CI] \\
                    \midrule
                    Fine-grained (4 classes) & 0.46 [0.37, 0.55] \\
                    Ambiguous vs Definitive  & 0.44 [0.35, 0.54] \\
                    Task vs Exo              & 0.78 [0.57, 0.93] \\
                    \bottomrule
                    \end{tabular}
                \end{table}

        \subsubsection{Comparison of prototype vocabularies}  

            In order to measure the extent to which different annotators' judgement yields different prototype vocabularies, we compare \textsc{Model-A} and \textsc{Model-B} in terms of (i) their cardinality, as an indicator of sparsity, (ii) their silhouette score, as a normative unsupervised measure of within-prototype consistency and across-prototype separability; (iii) their amount of accepted prototype, following the procedure described in \Cref{sec:chapter_5_results_consolidation}, and (iv) their semantic coverage, assessed by the overlap and differences between \textit{accepted} vocabulary items. 

            For the latter, prototypes in the larger \texttt{Model-A} vocabulary are matched to their most co-occurring prototype in the other vocabulary, based on embedding-level overlap. Overlapping and unmatched prototypes are recorded, providing an assessment of redundancy and coverage gaps between vocabularies (\Cref{sec:chapter_5_results_human_feedback}). %Lastly, we report an empirical taxonomy of the different scenarios observed, highlighting current weaknessess of our prototyping procedure and vocabulary consolidation stages, and sheding light on the challenges and perpectives of egocentric vision prototyping. % It should shed ligth on the different cases 
        
            % we perform Hungarian matching \cite{kuhn1955hungarian} between the two sets of prototypes, using the cosine distance between prototypes as a cost function. This yields a one-to-one mapping between prototypes of \textsc{Model-A} and \textsc{Model-B}, which we use to identify common and specific prototypes.

            \paragraph{Results.\\}
            %To assess the similarity of the prototype vocabularies generated by the two annotators, we computed the Jaccard similarity coefficient between the sets of accepted prototypes. This metric quantifies the overlap between the two sets, providing insight into the consistency of the visual semantics captured by each annotator. Our analysis revealed a Jaccard similarity of $0.65$, indicating a moderate level of agreement in the selected prototypes. Notably, the shared prototypes predominantly corresponded to visually salient actions, suggesting that both annotators were attuned to the same key aspects of the task.

            % (143237, 142727) vs  (332150, 417733)
            We observed that both vocabularies capture a similar range of egocentric actions, despite some differences in subjective judgments during centroids selection. We report in \Cref{fig:chapter_5_assignment_cost} the embedding-level assignment cost matrix, illustrating the matching each prototype of the larger \texttt{Model-A} vocabulary to their most co-occurring prototype in \texttt{Model-B}. 

            \begin{figure}[!htbp]
                \centering
                \includegraphics[width=1.0\linewidth]{chapters/chapter 5/figures/chapter_5_assignement_cost.png}
                \caption{\textbf{Embedding-level assignment cost matrix (row-normalized) between the two prototype vocabularies.} Each row represents a \textit{valid} prototype from \texttt{Model-A}, and each column a \textit{valid} prototype from \texttt{Model-B}. Columns are ordered using the Hungarian matching algorithm to maximize the alignment of matching prototypes along the diagonal. Darker diagonal structures indicate sharp one-to-one correspondences, while spread-out rows highlight ambiguous or redundant prototypes.}
                \label{fig:chapter_5_assignment_cost}
            \end{figure}

            The dominant diagonal structure indicates a sharp one-to-one correspondences, with a few spread-out rows revealing the redundant prototypes in both vocabularies. To verify this, we inspect the top-2 matches of all prototypes to confirm that they correspond to similar actions, as in the example of \Cref{fig:chapter_5_matches}. Additional examples are provided in \Cref{fig:chapter_5_matches_second_best} and \Cref{fig:chapter_5_best_match_recipe}. Importantly, we did not find prototypes that were not present in the other vocabulary, suggesting that both vocabularies effectively captured a similar range of egocentric actions. Second and third-best matches were occasionally encoding different action seemingly happening before or after the anchor action, which is identified as a minor cause of misalignment owing to transition phases between actions, inherently more ambiguous and revealing spurious differences between vocabularies. %We quantify next the global agreement and temporal alignment between the two vocabularies.


            \begin{figure}[!htbp]
                \centering
                \includegraphics[width=0.9\linewidth]{chapters/chapter 5/figures/chapter_5_matches.png}
                \caption{\textbf{Example of top-2 prototype matches based on the number of overlapping embeddings assigned to these prototypes.} The anchor prototype from \texttt{Model-A} (left) is matched to its best (middle) and second-best (right) counterparts from \texttt{Model-B}. In this example, all three prototypes consistently capture the action of "merging the dough", demonstrating semantic alignment between the two vocabularies while also illustrating their residual redundancies.}
                \label{fig:chapter_5_matches}
            \end{figure}


            %The mean argmax value of the matrix was $0.85 \pm 0.05$, indicating that most prototypes had a clear best match, while a mean split proportion of $0.15$ (threshold = $0.2$) and average entropy of $0.3$ highlighted occasional one-to-many mappings. %The row-wise normalized entropy of the assignment cost matrix averaged 0.409 (95\% CI: 0.347–0.471), indicating that prototype assignments were generally sharp, with moderate ambiguity for some cases.
        
        \subsubsection{Comparison on produced symbolic representations}

            % In order to evaluate the effect of using \textsc{Model-A} or \textsc{Model-B} on the participant's symbolic representation, we compute several unsupervised indicators of their quality and alignment: (i) their respective total assignment distance sum, as a proxy for data fitting, (ii) their number of estimates segments, and (iii) the Adjusted Rand index between symbolic sequences, as well as the amount of "pseudo-labels" and label switch they yielded. 
            Beyond the direct comparison of vocabularies, we examine their downstream impact on symbolic representations of participants’ egocentric vision trajectories. Specifically, we compute a set of unsupervised indicators, which together capture data fitting and temporal alignment. 

            \begin{itemize}

                \item \textbf{Adjusted Rand Index (ARI) \cite{Hubert1985}.}  
                The ARI, defined in \Cref{apd:chapter_5_clustering_metrics_definitions}, quantifies the similarity between two symbolic representations of the same participant by considering all pairs of estimated sample and counting pairs that are assigned in the same or different prototypes, correcting for lucky agreements. The ARI ranges from \(-1\) (indicating worse than random agreement) to \(1\) (perfect agreement), with \(0\) indicating agreement equivalent to random chance.

                \item \textbf{Normalized Mutual Information (NMI) \cite{Lancichinetti2009}.}

                The NMI, defined in \Cref{apd:chapter_5_clustering_metrics_definitions}, measures the mutual dependence between the two sets of prototypes, normalized by their entropies. A value of 1 indicates perfect correspondence, while 0 indicates no relationship.

                \item \textbf{Data fitting assessment.}  
                The total sum of cosine distances between each empirical embedding $x_i$ and its assigned prototype centroid $\mu_{c_i}^{[K]}$:  
                \[
                    d_\text{tot} = \sum_{i=1}^N d(x_i, \mu^{[K]}_{c_i}), 
                    \quad c_i = \arg \min_{k \in [K]} d(x_i, \mu^{[K]}_k).
                \]  
                Lower $d_\text{tot}$ values indicate embeddings are more tightly clustered around centroids, suggesting better representational alignment. Higher values may signal underfitting or inadequate granularity.  

                \item \textbf{Background assignment consistency.}  
                We compute both the proportion of frames labeled as background and their temporal correspondence across models, as measured by the ARI and NMI. This quantifies the degree to which both prototype vocabularies capture similar regions of low assignment confidence.

                \item \textbf{Assignment stability.}  
                The proportion of label switches between pointwise assignments and their segment-level consensus, $p^{\text{labels}}_{\text{change}}$, characterizes (i) the temporal regularization imposed by boundary detection, (ii) discrepancies in granularity between prototyping and segmentation, and (iii) assignment volatility. 

                \item \textbf{Pseudo change-points.}  
                We compute the proportion of change-points occurring between two identical consecutive labels, $p^{\text{withdrawn}}_{\text{CP}}$, capturing mismatches between the granularities of the multiple change-point detection and the prototypes temporal span. Frequent pseudo change-points suggest prototypes cover broader temporal scales than the segmentation, whereas rare pseudo change-points indicate finer-grained prototyping.% (\Cref{sec:chapter_5_comparison_across_models}). 

                \item \textbf{Temporal granularity.}
                We compare the empirical cumulative density function (ECDF) of effective action duration across models, that is, after the segment-level majority voting and the pseudo change-points withdrawal. For a set of sample $\{X_1, \dots, X_n\}$, the ECDF is defined as \[ \hat{F}_n(x) = \frac{1}{n} \sum_{i=1}^n \mathbf{1}_{\{X_i \leq x\}},\] where $\mathbf{1}_{\{\cdot\}}$ is the indicator function. In particular, the total number of estimated segments, $N_\text{segments}$, is used as a proxy indicator of overall temporal resolution: fewer segments reflect coarser groupings of actions, whereas more segments indicate finer partitioning. Non-informative background segments are withdrawn for this analysis.

            \end{itemize}

            Altogether, these indicators provide a comprehensive assessment of how prototype vocabularies derived from different annotators influence symbolic representations performances


            \paragraph{Results: comparison of non-background segments.\\}

                The Normalized Mutual Information (NMI) was 0.791 (95\% CI: [0.790, 0.792]) and the Adjusted Rand Index (ARI) was 0.682 (95\% CI: [0.680, 0.685]), indicating strong but imperfect agreement between the vocabularies. While the previous analysis of prototypes matches indicates that this elevated level of agreements reflects meaningful semantic alignment, the misalignments between symbolic representations seem mainly caused by prototypes redundancies within each vocabulary, leading to one-to-many label mappings that lower the agreement scores.

                \begin{figure}[!htbp]
                    \centering
                    \includegraphics[width=0.9\linewidth]{chapters/chapter 5/figures/segments_duration_models_comparison.png}
                    \caption{\textbf{Empirical cumulative density function (ECDF) of segments durations.} The ECDF illustrates the distribution of segment durations across participants for both models, demonstrating they capture similar temporal dynamics granularity.}
                    \label{fig:chapter_5_segments_duration_model_comparison}
                \end{figure}
                
                The total assignment distance sums, pseudo change-points proportion, and label switches indicators were similar for both models (see \Cref{tab:chapter_5_prototype_pipeline}), suggesting comparable data fitting quality and temporal granularity of the resulting symbolic representations (\Cref{tab:chapter_5_prototype_pipeline}). 
                
                The empirical cumulative density functions (ECDF) of effective action duration were also nearly identical across models (\Cref{fig:chapter_5_segments_duration_model_comparison}), with median durations of $6.4$ seconds for \texttt{Model-A} and $6.6$ seconds for \texttt{Model-B}. The total number of estimated segments was $N_\text{segments} = 1{,}247 \pm 290$ for \texttt{Model-A} and $N_\text{segments} = 1{,}298 \pm 310$ for \texttt{Model-B}, further confirming similar temporal resolutions.

            
            \paragraph{Results: comparison of background segments classification.\\}


                The background assignment confusion matrix (\Cref{tab:bg_confusion}) shows that both vocabularies agree on the vast majority of embeddings, with $80\%$ consistently classified as non-background and $12\%$ as background. While discrepancies remain limited, with $8\%$ of embeddings being assigned background by only one of the two vocabularies, suggesting different specialization of the model during the accumulation process. The ARI (mean = 0.627, 95\% CI = [0.625, 0.630]) and NMI (mean = 0.431, 95\% CI = [0.428, 0.434]) indicate a moderate level of concordance. This suggests that, despite general overlap in background detection, the vocabularies may differ in identifying low-confidence assignment.

                \begin{table}[!htbp]
                    \centering
                    \caption{\textbf{Confusion matrix of background assignments.} 
                    Agreement between \textsc{Model-A} (rows) and \textsc{Model-B} (columns) in assigning frames as background (bg) or non-background (non-bg).}
                    \label{tab:bg_confusion}
                    \vspace{0.5em}
                    \begin{tabular}{lcc}        
                        \toprule
                        & \textbf{B: Non-background} & \textbf{B: Background} \\
                        \midrule
                        \textbf{A: Non-background} & 0.80 & 0.04 \\
                        \textbf{A: Background}     & 0.04 & 0.12 \\
                        \bottomrule
                    \end{tabular}
                \end{table}

\section{Discussions}
\label{sec:chapter_5_conclusion}
    % Build it as benefits and then drwabacks and then future work ? 

    Our objective is to facilitate the study of disease-induced behavioral alterations by encoding the high-dimensional participants' egocentric vision trajectories into a lower-dimensional sequence of indexes, each representing interpretable registered prototypes. 
    %These are called symbolic representations in this work, and must accurately reflect the participant’s behavior to legitimate subsequent analysis and derive meaningful clinical interpretations. 
    We introduced a novel unsupervised, data-driven framework to derive fine-grained symbolic representations of egocentric videos. Our approach leverages the VideoMAE-v2 encoder together with an iterative human-in-the-loop refinement procedure to construct an interpretable vocabulary of egocentric vision prototypes. In contrast with body landmarks-based approaches which operates on spatial coordinates spaces, our prototypes encode participant's visual experience explicitly, capturing environmental cues accurately and promising richer but also more complex downstream analysis of human behavior. To enable fair comparison across symbolic representations in downstream analyses, we enforced a shared vocabulary with a consistent temporal granularity of actions, implemented through a global penalization term in the change-point detection objective. Importantly, while supervision was used to monitor clustering estimates, the definition of actions remained fully unsupervised. Finally, two independent annotators produced vocabularies that, despite minor subjective differences in centroid selection, exhibited strong semantic alignment and temporal consistency in the resulting symbolic representations. 

    Our results demonstrate that the estimated prototypes capture a broad range of actions with strong generalizability to unseen participants. This generalizability, validated empirically by cross-validation, is theoretically attributable to the \textit{prescribed} nature of the task performed in reproducible setting, and the consistent use of the same head-mounted camera and video encoder. We observed that several prototypes encode similar actions due to two main factors: (i) the different strategies adopted by participants to execute the task, and (ii) variations in ingredient packaging and materials, leading to differences in the egocentric video recordings. 
    The former highlights the importance of vocabulary granularity, i.e., the scale at which prototypes should be grouped depends on the analysis. For instance, breaking the eggs may occur using the trash, the sink, or the bowl at participant's disposal. While these three visual atoms could be merged for action-level analysis, they may also remain distinct to explicitly encode alternative strategies. The latter underscores the importance of maximizing visual reproducibility in experimental conditions (e.g., fixed material and ingredients) to minimize spurious prototype ambiguity and facilitate downstream analysis.

    The hierarchical clustering procedure, while effective at reducing redundancy, occasionally merges centroids corresponding to distinct actions, resulting in semantically ambiguous prototypes. This limitation suggests that the structure induced by the distribution of pairwise centroids similarities constrains the recovery of mutually-exclusive prototypes. Specifically, representations of semantically distinct actions in the latent space can be closer than slightly different instances of semantically similar atomic actions, making it difficult for clustering algorithms to disentangle them effectively. This highlights the importance of the manual inspection step in ensuring the semantic integrity of the final vocabulary. Future work could explore more advanced consolidation strategies, for example by integrating complementary sources of information such as temporal dynamics, to better capture the subtleties of egocentric actions. For the behavioral analysis of the next chapter, we maximize the accuracy of the final vocabulary by combining \textit{valid} prototypes of the $2^{nd}$ consolidation, prototypes of the $1^{st}$ consolidation for \textit{ambiguous} cases, and the original validated centroids in case these were also classified as \textit{ambiguous}.

    To conclude, we see this work as a promising proof of concept for the computation of fine-grained symbolic representation of egocentric video for behavioral analysis. As discussed in the next chapter, these interpretable symbolic representation offer rich material to incorporate and advance the integration of multimodal data, including eye-tracking, inertial measurement units (IMU), and audio recordings, to comprehensively characterize participant's behavior in naturalistic settings.


    \paragraph{Future directions of research.\\}

        Several avenues could extend this work:
        \begin{itemize}
            \item \textbf{Refined definition of the residual space.} The current residual space definition yields redundant accumulated centroids. Refined discrimination of misrepresented samples could enhance novelty detection and accelerate convergence rate of the procedure. 
            \item \textbf{Temporal information integration.} Incorporating temporal cues to better discriminate visually similar yet semantically distinct actions. This could be achieved using early fusion of timestamps into the centroid accumulation process, for instance via temporal positional embeddings.
            \item \textbf{Task-specific fine-tuning.} Adapting the video encoder to the Smartflat task to improve samples discrimination and facilitate prototyping. 
            \item \textbf{Automating human feedback.} Reducing the cost of annotation by leveraging the existing bank of annotated centroids to train a classifier capable of predicting the validity of newly estimated centroids.

        \end{itemize}


    